\documentclass[11pt]{article}
\usepackage[T1]{fontenc}
\usepackage[utf8]{inputenc}
\usepackage{mathptmx}
\usepackage[a4paper,margin=2cm]{geometry}
\usepackage{amsmath,amssymb}
\usepackage{url}
\usepackage[hidelinks]{hyperref}
\usepackage{natbib}

\title{Explored literature}
\date{10 October 2026}

\begin{document}
\maketitle

This document records what was read and is not an L row of
\path{MANUSCRIPT.tex}. Each record under \path{lit/} holds the quotations,
their pages and the Lean that checks our reading. Locators are to the version
read.

\section{\citet{bayraktar2026}}

Read in full, arXiv:2603.14557v2 (record \path{lit/bayraktar_2026/}).

\paragraph{The model of this bundle.} The cap, the diffusion of the
asset-to-deposit ratio and the switching point between the solvency and the
liquidity branch are this paper's (eq.~4, p.~10, p.~15). L1 rests on it.

\paragraph{The boundary at one.} Remark~3.1 (p.~11) states that at the
zero-equity level the cap forces zero risky exposure, the diffusion vanishes,
and the point is locally repelling when $r>r_L$ and attracting when $r<r_L$.
The paper avoids the point by intervening at an exogenous threshold above one
(p.~2). This bears on H2 and H2b of \path{TODO.md}, whose degeneracy and
repelling condition are therefore in print. What the paper does not do is
change coordinate, so the reading of the degeneracy as an artefact of the
level coordinate is not in it.

\paragraph{No fixed issuance cost.} Issuance costs are a dilution cost and a
proportional cost. The paper states that a fixed issuance cost would
generally break the one-dimensional reduction (p.~9) and that it could make
earlier recapitalisation or inaction at the boundary optimal (p.~13). A fixed
cost keeps the reduction only when it is charged per unit of liabilities, as
in \path{PROOFS_v2.tex}. This bears on H3 and H3b.

\paragraph{Smooth fit.} Smooth fit is proved at the dividend barrier, and
issuance occurs only at the exogenous threshold (Theorem~3.1, p.~12). There
is no smooth fit at an optimally chosen recapitalisation trigger. This bears
on H4.

\section{\citet{barberis1999}}

Read in full in NBER Working Paper 7220 (July 1999). The manuscript cites
the journal version, \citet{barberis2001}, read in full as well. Record
\path{lit/barberis_huang_santos_1999/}.

\paragraph{Kinked-linear loss aversion.} The loss-aversion term is piecewise
linear and kinked at the origin (p.~10), and the paper drops the curvature of
prospect theory for simplicity (p.~11). A kinked-linear penalty is therefore
not new as such. This bears on H5, whose contrast with S-shaped penalties
must not suggest otherwise.

\paragraph{Changing risk aversion.} Volatility above that of consumption
growth comes from time variation in effective loss aversion through the
benchmark state (Section~3). The level of loss aversion alone fails on both
volatility and the premium (p.~15).

\section{\citet{barberis2001}}

Read in full in the journal version (record
\path{lit/barberis_huang_santos_2001/}). L2 cites it. The journal version adds a second economy in which dividends and
consumption are separate, moves the constant-loss-aversion model to
Section~VI, where Table~XIII shows return volatility of 12 percent at every
$b_0$ (p.~47), and changes the calibration. The record lists the differences.

\section{\citet{li2023}}

Read in full, arXiv:2108.02648v4 (record \path{lit/li_yu_zhang_2023/}). L4
rests on it. The paper is also the S-shaped case that needs a concave
envelope (p.~1), against which the concave kinked-linear penalty of H5 is
set. Its Remark~4.1 attributes to the loss-aversion degree a dependence that
its own limit formulas carry only when the curvatures on gains and losses are
equal (pp.~33 to 34).

\section{\citet{engle2018}}

Read in full in the published article (record
\path{lit/engle_siriwardane_2018/}).

\paragraph{Precautionary capital.} The equity a firm must raise today to meet
a capital ratio in a crisis with a given confidence (Section~3.1). It is a
regulatory quantity of a firm, not of a control problem.

\paragraph{The leverage multiplier.} Equity volatility is approximately the
leverage multiplier times asset volatility (eq.~5, Appendix~A). The
multiplier rises with leverage and is around 3 for the median firm
(Section~2.3).

\bibliographystyle{plainnat}
\bibliography{refs}

\end{document}
